\section{A model-rocket flight, from first principles}\label{sec:primer}

This chapter is for the engineer who has never flown a model rocket. It builds, from first
principles, every physical concept the rest of the document leans on. Pure theory is set
apart in \emph{Physical background} boxes so it is never conflated with claims about the
code; every statement about what QtRocket actually does carries a source citation into the
pinned commit \code{254cd41}, and the chapter closes with a table mapping each concept to
the class that implements it and the chapter that examines it.

\subsection{Anatomy of a rocket, anatomy of a flight}\label{sec:primer-anatomy}

A sport model rocket is a slender axisymmetric stack: a \emph{nose cone} at the top, one or
more cylindrical \emph{body tubes}, a set of three or four \emph{fins} at the aft end, and a
replaceable solid-propellant \emph{motor} in a mount at the tail. Real airframes also carry a
recovery device (parachute or streamer) packed in the body tube and a launch lug or rail
button that guides the rocket along a rod or rail for its first meter of flight, before the
fins have enough airspeed to work. QtRocket's part catalog covers the structural core of this anatomy
--- nose cone, body tube, fin set, plus a hollow sphere and the motor leaf --- assembled
into an ownership tree (\cref{sec:parts}); recovery devices and launch guides are among the
part types that do not exist yet (\cref{sec:incomplete}).

A flight divides into four phases punctuated by one instant, and the vocabulary recurs
throughout this document (\cref{fig:primer-flight-phases}):

\begin{description}
\item[Pad.] The rocket sits on the launch rod at altitude zero until the igniter fires.
  QtRocket models the pad as a kinematic clamp: until the first time step that ends above
  $z = 0$ under power, position and velocity are held at rest, standing in for the normal
  force a real pad exerts \src{sim/Propagator.cpp}{105}. A rocket that cannot lift off is
  detected --- more than $3\unit{s}$ without exceeding $1\unit{m}$ --- and the run ends with a
  typed \code{NoLiftoff} verdict rather than silently \src{sim/Propagator.cpp}{146},
  \src{sim/Propagator.h}{26}.
\item[Boost.] The motor burns, typically for half a second to a few seconds, producing a
  thrust force that varies strongly over the burn. Production code ignites the motor at
  exactly $t = 0$ \src{model/RocketModel.cpp}{108}. During boost the rocket both
  accelerates and \emph{lightens}: the propellant leaves through the nozzle.
\item[Coast.] After burnout the rocket is a projectile decelerating under gravity and drag.
  In a real motor a slow-burning \emph{delay charge} smolders through this phase; QtRocket's
  motor model reports zero thrust after the last thrust sample \src{model/ThrustCurve.cpp}{51}
  and the vehicle simply coasts.
\item[Apogee.] The instant of peak altitude, where vertical velocity passes through zero.
  This is where a real flight fires the ejection charge and deploys recovery. QtRocket
  records apogee as a running maximum folded into the flight statistics every step
  \src{sim/TrajectoryStatistics.h}{26}; nothing else happens there.
\item[Descent.] A real rocket descends under parachute at a few meters per second. QtRocket
  has no recovery events at the pinned commit: every flight descends \emph{ballistically},
  under the same drag law as ascent, and terminates when the model reports that it is both
  below the launch site and descending --- $z < 0$ and $v_z < 0$
  \src{model/RocketModel.cpp}{65}. Ejection-delay data is parsed by the local motor-file
  loaders and persisted, but no flight logic consumes it \src{model/MotorModel.h}{303}. The consolidated
  inventory of unbuilt flight events lives in \cref{sec:incomplete}.
\end{description}

\tikzfig{flight-phases}{The phases of a model-rocket flight, annotated with the code that
observes each one. QtRocket holds the vehicle on the pad until the first powered step ends
above $z=0$ \srccap{sim/Propagator.cpp}{105}, ignites the motor at $t=0$
\srccap{model/RocketModel.cpp}{108}, tracks apogee as a running maximum
\srccap{sim/TrajectoryStatistics.h}{26}, and --- because no recovery events exist at the
pinned commit --- flies the descent ballistically until the termination predicate $z<0$ and
$v_z<0$ fires \srccap{model/RocketModel.cpp}{65}.}{fig:primer-flight-phases}

\subsection{The three forces}\label{sec:primer-forces}

Three forces act on a rocket in flight, and QtRocket's production force model assembles
exactly these three terms and no others \src{model/RocketModel.cpp}{68}.

\paragraph{Thrust.} A solid motor's thrust is not a constant; it is a measured function of
time $F(t)$, published as a \emph{thrust curve} --- a table of (time, force) samples from a
certification test stand. A typical curve spikes to its peak in the first tenth of a second
(high initial acceleration gets the rocket moving fast enough for fin stabilization), then
settles to a sustainer plateau, then tails off at burnout. QtRocket stores the curve as
sampled pairs and interpolates linearly between samples, returning zero outside the burn
window \src{model/ThrustCurve.cpp}{47}. In flight, the thrust force is applied along the
world vertical and assumed to act through the center of mass
\src{model/RocketModel.cpp}{71}; since no attitude state exists in the dynamics
(\cref{sec:primer-ode}), the rocket's axis and the world $z$-axis are the same thing.

\paragraph{Weight.} Weight is $m(t)\,g$ --- and the time argument matters. The propellant
that produces the thrust leaves the vehicle as it burns, so the mass in Newton's second law
falls throughout boost: propellant can be up to half of a small motor's loaded mass. QtRocket obtains $m(t)$ from the part tree's composite mass, with the motor as a
time-aware leaf part
(\cref{sec:parts,sec:motors}), and evaluates gravity through a pluggable model --- uniform
$-g_0\hat{z}$ or Newtonian inverse-square --- at the integrator's trial position
\src{model/RocketModel.cpp}{77} (\cref{sec:environment}).

\paragraph{Drag.} Air resistance is the force a newcomer is most likely to underestimate:
for a lightweight airframe at a hundred meters per second it can rival or exceed the
rocket's weight, and it is the reason drag-free vacuum apogees run to multiples of real
ones.

\begin{physbox}[The quadratic drag law, symbol by symbol]
For a blunt body moving faster than a slow walk, drag is well modeled by
\begin{equation}
  D = \tfrac{1}{2}\,\rho\,v^{2}\,C_d\,A .
\end{equation}
Each symbol earns its place:
\begin{itemize}
\item $\rho$ is the local \emph{air density} (kg/m$^3$). Drag is momentum transferred to
  the air the vehicle displaces, so it scales with how much air mass occupies each swept
  meter; density falls with altitude (\cref{sec:primer-atmosphere}).
\item $v^{2}$: the vehicle sweeps through air proportionally to $v$ and imparts to that air
  a velocity change also proportional to $v$ --- hence the square, which makes the equation
  of motion nonlinear (\cref{sec:primer-integration}).
\item $C_d$ is the dimensionless \emph{drag coefficient}, the shape's efficiency at evading
  the momentum bill: about 1.2 for a flat disc, 0.3--0.75 for typical model rockets.
\item $A$ is the \emph{reference area} (m$^2$) --- a bookkeeping convention, not
  independent physics: $C_d$ is only meaningful relative to a stated area, and the hobby
  convention is the airframe's maximum frontal cross-section (the widest body disc, never
  inflated by fin span).
\end{itemize}
A useful corollary is \emph{terminal velocity}, the speed at which drag balances weight in
steady descent: $v_t = \sqrt{2mg/(\rho C_d A)}$.
\end{physbox}

QtRocket implements the vector form of this law in one line, written with $|v|\,\vec v$
rather than $v^2\hat v$ so that zero velocity produces exactly zero drag with no division
\src{model/RocketModel.cpp}{88}:
\begin{equation}
  \vec F_{\text{drag}} \;=\; -\tfrac{1}{2}\,\rho(z)\;\lVert\vec v\rVert\;
  \mathtt{dragCoefficient}\cdot\mathtt{referenceArea}\;\vec v ,
\end{equation}
where $\rho(z)$ is sampled from the selected atmosphere model at the trial altitude, clamped
to $z \ge 0$ \src{model/RocketModel.cpp}{85}. Both coefficients are user-set scalars with
defaults $C_d = 1.0$ \src{model/RocketModel.h}{138} and $A = 1.134\times10^{-3}\unit{m^2}$
--- the frontal disc of a 38\,mm body tube \src{model/RocketModel.h}{141}. A helper that
derives the reference area from the assembled geometry under the max-frontal-disc convention
exists \src{model/RocketModel.cpp}{48}, but its wiring status belongs to \cref{sec:aero}.
The drag term is also the best-tested physics in the tree: the terminal-velocity force
balance above is pinned analytically by an integration test
\src{tests/PhysicsIntegrationTests.cpp}{262}, and drag-versus-vacuum apogee sweeps run across
every reference airframe and the full motor ladder \src{tests/DesignMatrixTests.cpp}{515}.

\tikzfig{freebody-stability}{Left: the free-body diagram of the point-mass flight model ---
thrust $F(t)$ along the world vertical \srccap{model/RocketModel.cpp}{71}, weight $m(t)g$
\srccap{model/RocketModel.cpp}{77}, and quadratic drag opposing the velocity vector
\srccap{model/RocketModel.cpp}{88}. Right: the geometry of static stability --- when a gust
tilts the rocket to a small angle of attack, the fins' normal force acts at the center of
pressure (CP); with the CP aft of the center of gravity (CG) the resulting moment restores
the rocket toward zero angle, and the CP--CG gap measured in body diameters is the static
margin. QtRocket computes both stations on a shared datum (\srccap{sim/Aero.h}{25}) but no
torque acts in flight at the pinned commit (\srccap{model/RocketModel.cpp}{94});
\cref{sec:aero} carries the analysis.}{fig:primer-freebody}

\subsection{Newton's second law as the equation the simulator integrates}%
\label{sec:primer-ode}

Summing the three forces gives the equation of motion. Because the mass changes, it is worth
writing carefully:
\begin{equation}
  \dot{\vec x} = \vec v, \qquad
  \dot{\vec v} = \frac{\vec F_{\text{thrust}}(t) + \vec F_{\text{gravity}}(\vec x, t)
                      + \vec F_{\text{drag}}(\vec x, \vec v)}{m(t)} .
  \label{eq:primer-ode}
\end{equation}
The second-order law $\ddot{\vec x} = \vec F/m$ has been reduced to a coupled first-order
pair --- the standard trick that lets a general-purpose ODE (ordinary differential equation)
solver advance position and velocity together. The pair
$(\vec x, \vec v) \in \mathbb{R}^6$ is the \emph{state vector}: the minimal set of numbers
from which the future of the system is computable. QtRocket builds
\cref{eq:primer-ode} verbatim, as a lambda constructed once in the propagator's constructor
(\cref{lst:primer-ode}); this closure \emph{is} the entire integrated physics, and
\cref{sec:simcore} dissects the loop that drives it.

\begin{listing}[htbp]
\begin{minted}{cpp}
// Linear ODEs: dx/dt = v, dv/dt = F/m. t is passed through so getForces() can evaluate a
// time-dependent forcing function (motor thrust curve) at the actual simulation time.
std::function<std::pair<Vector3, Vector3>(double, Vector3&, Vector3&)> linearODEs = [this](double t, Vector3& state, Vector3& rate) -> std::pair<Vector3, Vector3>
{
    Vector3 dPosition;
    Vector3 dVelocity;
    dPosition = rate;
    dVelocity = object->getForces(t, state, rate, *environment) / object->getMass(t);

    return std::make_pair(dPosition, dVelocity);
};
\end{minted}
\caption{Newton's second law with time-varying mass, as QtRocket states it
\srccap{sim/Propagator.cpp}{36}. The time argument is threaded through so the thrust curve
and the burning mass are sampled at each solver stage's node time rather than frozen at the
step start.}\label{lst:primer-ode}
\end{listing}

\begin{keyinsight}[Where the rocket equation went]
Textbook treatments of variable-mass motion add a momentum-flux term $\dot m\,v_e$ (exhaust
mass flow times exhaust velocity) --- the term that becomes the Tsiolkovsky rocket equation.
QtRocket carries no such term, and that is correct, not an omission: the thrust curve is a
\emph{measured force} from a test stand, so the momentum flux is already embedded in $F(t)$.
Mass loss enters the dynamics only through the divisor $m(t)$
\srccap{sim/Propagator.cpp}{41} and the weight term \srccap{model/RocketModel.cpp}{77};
adding the flux on top of a measured curve would double-count it.
\end{keyinsight}

\paragraph{Degrees of freedom, precisely.} A \emph{3-DOF (degree-of-freedom) point-mass}
model integrates the three translational coordinates only: the state is
$(\vec x, \vec v) \in \mathbb{R}^6$, forces are applied at a single point, and the vehicle
has no orientation --- it cannot pitch, yaw, roll, fly at an angle of attack, or respond to
a torque. A \emph{6-DOF rigid-body} model adds the three rotational degrees of freedom: an
attitude (conventionally a unit quaternion) and an angular velocity, giving a
thirteen-element state, with torques and the time-varying inertia tensor entering through
Euler's rigid-body equations. This is the single most important fidelity boundary in flight
simulation: stability, weathercocking, gravity turns, and spin all live on the rotational
side.

QtRocket today integrates exactly the three linear degrees of freedom
(\cref{lst:primer-ode}). The rotational half exists as deliberate, typed scaffolding with
no writers: the state record reserves quaternion, direction-cosine-matrix, and Euler-angle
slots that no code writes \src{sim/StateData.h}{46}; the model's torque accessor returns the
zero vector and has no caller in the simulation loop \src{model/RocketModel.cpp}{94}; and
the orientation integrator is a commented-out member with a design note
\src{sim/Propagator.h}{103}. One rotational input \emph{is} live: composite mass, center of
gravity, and inertia tensor are recorded into every state sample, expressly as ``the seam
the 6-DOF rotational ODE will read'' \src{sim/StateData.h}{59}. \Cref{sec:simcore} covers
the integrated path; \cref{sec:incomplete} inventories the dormant half.

\subsection{Stability: the center of gravity, the center of pressure, and the
arrow}\label{sec:primer-stability}

Nothing about a finless tube makes it fly straight. Point-mass simulation sidesteps the
question by construction; understanding what it sidesteps requires two points on the
airframe.

\begin{physbox}[Static stability and the Barrowman method]
The \emph{center of gravity} (CG) is where the vehicle's mass balances; the rocket rotates
about it. The \emph{center of pressure} (CP) is where the net aerodynamic \emph{normal}
force --- the sideways lift that appears when the airflow meets the body at a small
\emph{angle of attack} $\alpha$ --- effectively acts. The classic intuition is the arrow:
feathers at the back add lifting area far aft, dragging the CP behind the CG, so any
disturbance that tilts the arrow produces a force behind the pivot that weathervanes it back
into the wind. Fins are the rocket's feathers. If the CP sits \emph{ahead} of the CG, the
same moment amplifies the disturbance and the rocket tumbles.

The gap is quoted in \emph{calibers} --- body diameters:
$\mathrm{SM} = (x_{CP} - x_{CG})/d$ with stations measured tip-aft-positive; the hobby rule
of thumb is one to two calibers of positive margin. The same mechanism has a cost:
\emph{weathercocking}. A stable rocket leaving the rod in a crosswind flies at a momentary
angle of attack, and its own restoring moment turns it \emph{into} the wind, trading
altitude for upwind travel.

Computing the CP from geometry alone was solved for hobby rockets by James and Judith
Barrowman (1966). Their method linearizes each component's normal force at small $\alpha$
--- only the slope $C_{N_\alpha}$ per radian matters --- and applies slender-body theory: a
constant-diameter tube contributes zero normal-force slope; any nose cone contributes
exactly $C_{N_\alpha} = 2$ referenced to its base area, with its CP at a fixed fraction of
its length; fins contribute a closed-form term with a body-interference factor. Because all
components share one $\alpha$ and one reference area, the slopes add, and the composite CP
is the slope-weighted average of the component CPs. The result is a closed-form,
microsecond-cheap calculation that is accurate in exactly the regime sport rockets occupy
--- slender, axisymmetric, small $\alpha$ --- which is why essentially every hobby
simulator, OpenRocket included, is built on it.
\end{physbox}

QtRocket implements the Barrowman method faithfully at the component level --- the cone's
$C_{N_\alpha} = 2$ at base area \src{model/parts/ConicalNoseCone.cpp}{64}, the $N$-fin term
with the interference factor \src{model/parts/FinSet.cpp}{77} --- and composes the pieces
onto the same datum as the composite CG so that the static margin is literally a
subtraction, \code{cp()}~$-$~\code{cg()} \src{sim/Aero.h}{25},
\src{model/parts/Part.cpp}{231}. At the pinned commit this machinery is \emph{dormant}:
built and unit-tested with zero production consumers --- the test suite states ``nothing in
production consumes this'' \src{model/tests/AeroTests.cpp}{18} --- and since no torque acts
and no angle of attack is ever computed, no flight is affected by stability or its absence.
Weathercocking is doubly out of reach: it needs both the restoring moment and a wind field,
and the wind model is a zero-returning stub with no call sites \src{sim/WindModel.cpp}{16}.
\Cref{sec:aero} gives the Barrowman layer the full treatment, including a verified sign
defect in the fin-set frame that must be fixed before the seam is consumed.

\subsection{The atmosphere}\label{sec:primer-atmosphere}

Drag needs $\rho(z)$, and $\rho$ is not a constant: a rocket that climbs $3\unit{km}$ flies
through air about a quarter thinner than at the pad.

\begin{physbox}[Hydrostatic balance and the standard atmosphere]
Air density falls with altitude because each layer of the atmosphere must carry the weight
of all the air above it. A static column obeys the \emph{hydrostatic equation}
\begin{equation}
  \frac{dP}{dh} = -\rho\, g ,
\end{equation}
and the ideal gas law $P = \rho \frac{R^{*}}{M} T$ ties pressure, density, and temperature
together. Substituting one into the other shows pressure decaying roughly exponentially with
altitude, with a scale height near $8\unit{km}$ --- but the exact profile depends on the
temperature $T(h)$, which real weather varies daily. A \emph{standard atmosphere} resolves
this by decree: it prescribes a fixed, piecewise-linear temperature profile (the troposphere
cools at $6.5\unit{K/km}$; the stratosphere is first isothermal, then warms as ozone absorbs
ultraviolet), which makes the hydrostatic equation integrable in closed form layer by layer
--- an exponential in the isothermal layers, a power law in the gradient layers. The result
is a reproducible reference profile, the same for every user on every day, which is exactly
what a simulator wants.
\end{physbox}

QtRocket ships three pluggable atmospheres --- constant sea-level, vacuum, and the US
Standard Atmosphere 1976 --- selectable at runtime behind one interface, with constant
sea-level as the default \src{sim/Environment.h}{31}. The US Standard 1976 implementation
codes the seven layer tables \src{sim/USStandardAtmosphere.cpp}{21} and the closed-form
per-layer solutions \src{sim/USStandardAtmosphere.cpp}{93}, and is validated against the
published NOAA report tables at $0.1$--$0.5$\% relative tolerance from $0$ to $80\unit{km}$
\srcf{sim/tests/USStandardAtmosphereTests.cpp}. Of the five properties the atmosphere
interface exposes, the flight path consumes exactly one --- density
\src{model/RocketModel.cpp}{86}; temperature, pressure, speed of sound, and viscosity are
computed and await a Mach- or Reynolds-aware drag model. \Cref{sec:environment} covers the
environment layer, including gravity models and the geoid.

\subsection{Why numerical integration}\label{sec:primer-integration}

\Cref{eq:primer-ode} has no closed-form solution for any realistic flight. Three
ingredients each independently destroy solvability: the drag term is quadratic in $v$
(nonlinear), the thrust forcing $F(t)$ is an arbitrary tabulated curve (not an elementary
function), and both $m(t)$ and $\rho(z)$ vary along the trajectory. Even the drastically
simplified case --- vertical flight, constant thrust, constant mass, constant density ---
yields hyperbolic-tangent and logarithmic expressions that fall apart the moment thrust
varies. So every flight simulator marches the state forward numerically, and the craft lies
in doing so accurately and cheaply.

\begin{physbox}[From Euler to Runge--Kutta, and why steps adapt]
The simplest scheme is Euler's: sample the derivative at the step start and walk straight
along it, $x_{n+1} = x_n + h\,f(t_n, x_n)$. Its global error shrinks only linearly with the
step size $h$ --- halving the error costs double the work. The classical fourth-order
Runge--Kutta method (RK4) instead probes the derivative four times per step (at the start,
twice at midpoint estimates, and at the end) and takes a weighted average; the leading
error terms cancel, so global error scales as $h^4$: halving $h$ buys sixteen times the
accuracy. The remaining inefficiency is the \emph{fixed} step: a flight needs tiny steps
during the violent first tenth of a second of the burn and could take huge ones during a
placid coast. \emph{Embedded} methods such as Runge--Kutta--Fehlberg 4(5) compute a fourth-
and a fifth-order estimate from the same six derivative samples; their difference is a free
estimate of the local error, which a controller compares against a tolerance --- shrinking
the step where the physics is fast, growing it where it is smooth.
\end{physbox}

QtRocket provides both: a textbook fixed-step RK4 \src{sim/RK4Solver.h}{45} and an adaptive
Runge--Kutta--Fehlberg 4(5) with the classical tableau \src{sim/RK45Solver.h}{75}, at a
default tolerance of $10^{-6}$ that no production caller overrides \src{sim/RK45Solver.h}{158}, selectable by name
at runtime from either front end. The payoff of adaptivity is measured in the tree: the
same test flight that costs RK4 roughly 4\,050 steps costs the adaptive solver about 530
\src{tests/PhysicsIntegrationTests.cpp}{163}. Two details from the box above are
load-bearing in the code. First, every solver stage passes its own node time into the ODE,
so the thrust curve is sampled at the right instants within a step \src{sim/RK4Solver.h}{55}
--- freezing thrust at the step start once inflated adaptive-run apogees by roughly 8\%
\src{tests/PhysicsIntegrationTests.cpp}{149}. Second, adaptive step growth must be capped
\src{sim/RK45Solver.h}{144}: coasting in vacuum is integrated \emph{exactly} by both
embedded orders, the error estimate collapses to zero, and an uncapped controller would
grow the step without bound and leap past apogee and the ground. \Cref{sec:simcore}
examines both solvers, the run loop, and its typed termination taxonomy
\src{sim/Propagator.h}{35}.

\subsection{Motors: impulse classes, designations, and the thrust
curve}\label{sec:primer-motors}

\begin{physbox}[Reading a motor designation]
Hobby motors are certified and labeled by \emph{total impulse}, the area under the thrust
curve, $I_{\text{tot}} = \int F\,dt$ in newton-seconds. Impulse classes are letters, and
each full letter doubles the ceiling of the one before: class A tops out at
$2.5\unit{N\,s}$, B at $5$, C at $10$, and so on up through G ($160$, the top of the
``model rocket'' range), H--I (level-1 high power), J--L (level 2), and M ($10\,240$,
level 3); fractional classes 1/4A and 1/2A sit below A. A designation like \emph{C6-5}
reads: class C (total impulse in the 5--10$\unit{N\,s}$ bracket), average thrust
$6\unit{N}$, and a $5\unit{s}$ ejection delay between burnout and the recovery charge.
Manufacturer suffixes, such as the T in \emph{G80T}, name the propellant formulation. Two
motors of the same class can fly very differently: a high-thrust short burn and a gentle
long burn deliver the same impulse with different peak accelerations --- which is why
simulators consume the full thrust curve, never just the letter.
\end{physbox}

QtRocket's motor subsystem (\cref{sec:motors}) mirrors this taxonomy directly. The impulse
class is derived from the leading letter of the motor code, with the fractional classes
1/4A and 1/2A special-cased \src{model/MotorModel.cpp}{112}; the heavy test suite flies the
complete 1/4A$\to$M ladder across every reference airframe
\src{tests/DesignMatrixTests.cpp}{515}. Delay options are parsed from the local ingest
formats and stored --- with $1000$ as the sentinel for a plugged motor
\src{model/MotorModel.h}{303} --- though, as \cref{sec:primer-anatomy} noted, no flight
logic yet consumes them. The thrust curve is sampled with linear interpolation and careful
edge handling: curves whose first published sample sits at $t > 0$ ramp from an implicit
origin \src{model/ThrustCurve.cpp}{76}. And because $m(t)$ matters
(\cref{sec:primer-forces}), the motor also derives a burning-mass curve: a 128-point table
that depletes propellant in proportion to the fraction of total impulse delivered --- the
standard constant-specific-impulse assumption \src{model/MotorModel.cpp}{132},
\src{model/MotorModel.cpp}{141} --- so mass falls fastest exactly when the curve burns
hardest. The motor's three mass regimes --- loaded before ignition, depleting during the
burn, casing after burnout \src{model/MotorModel.cpp}{28} --- are the single time-varying
input the part tree's composite mass walk sees (\cref{sec:parts}).

\subsection{From concept to code}\label{sec:primer-map}

\Cref{tab:primer-map} closes the primer by mapping each physical concept introduced above
to the code that owns it and the chapter that examines it. The table doubles as a capability
summary at the pinned commit: the point-mass rows are complete and consumed; the rotational
and recovery rows name machinery that is built but dormant, or absent outright.

\begin{table}[htbp]
\centering
\small
\begin{tabularx}{\textwidth}{@{}>{\raggedright\arraybackslash}p{0.28\textwidth}%
                             >{\raggedright\arraybackslash}X l@{}}
\toprule
\textbf{Physical concept} & \textbf{Where QtRocket implements it} & \textbf{Chapter} \\
\midrule
Thrust curve $F(t)$ & \code{ThrustCurve} sampling + \code{MotorModel} burnout latch
  \srcf{model/ThrustCurve.cpp} & \cref{sec:motors} \\
\addlinespace
Burning mass $m(t)$ & 128-point impulse-proportional mass curve; \code{part::Motor} leaf
  in the tree \srcf{model/MotorModel.cpp} & \cref{sec:motors} \\
\addlinespace
Mass, CG, inertia of the airframe & \code{Part} composite tree, two-pass walk behind the
  mass-delta cache \srcf{model/parts/Part.cpp} & \cref{sec:parts} \\
\addlinespace
Geometry: where parts sit & \code{StationLink} intent + the placement resolver
  \srcf{model/parts/Placement.cpp} & \cref{sec:placement} \\
\addlinespace
Weight / gravity field & \code{ConstantGravityModel}, \code{SphericalGravityModel}
  \srcf{sim/GravityModel.h} & \cref{sec:environment} \\
\addlinespace
Air density $\rho(z)$ & \code{AtmosphericModel} registry; US Standard Atmosphere 1976
  \srcf{sim/USStandardAtmosphere.cpp} & \cref{sec:environment} \\
\addlinespace
Quadratic drag & Three-term force model in \code{RocketModel::getForces()}
  \srcf{model/RocketModel.cpp} & \cref{sec:aero} \\
\addlinespace
Newton's law, state, flight phases & \code{Propagator} run loop, pad hold, typed
  termination \srcf{sim/Propagator.cpp} & \cref{sec:simcore} \\
\addlinespace
Numerical integration & \code{RK4Solver} (fixed), \code{RK45Solver} (adaptive) behind
  \code{DESolver} \srcf{sim/RK45Solver.h} & \cref{sec:simcore} \\
\addlinespace
CP, static margin (Barrowman) & \code{Part::getCompositeAero()} onto the shared datum ---
  dormant \srcf{sim/Aero.h} & \cref{sec:aero} \\
\addlinespace
Attitude / 6-DOF, wind & Reserved seams: unwritten \code{StateData} slots, zero torques,
  stub \code{WindModel} & \cref{sec:incomplete} \\
\addlinespace
Recovery events, ejection delays & Delays parsed and persisted; no flight logic consumes
  them & \cref{sec:incomplete} \\
\bottomrule
\end{tabularx}
\caption{The primer's concepts mapped to their implementations at commit \code{254cd41}.
Every row's owning chapter carries the citations and the analysis; rows marked dormant or
absent are consolidated in \cref{sec:incomplete}.}\label{tab:primer-map}
\end{table}
